%=====================================================================
\section{K3 lattice realisation of the 140-label representation}\label{sec:rev8k3}
%=====================================================================
This section (revision~8, 19 October 2026) adds the results of [S38] (\path{2026-10-19-K3-140-check.md} and the folder \path{2026-10-19-Implications-checks/k3_140/}). Category convention as in \S\ref{sec:cat}: \textbf{Category~A} is proved, or exact (exact finite computation; a mod-$\ell$ computation is stated as such); \textbf{Category~A (cited)} means proved \emph{conditional only on} the published theorems named, which are not re-proved here; \textbf{Category~B} is floating-point evidence, citation-only identification, interpretation, analogy or conjecture. The two categories are kept separate and every statement is labelled.

\paragraph{Labels and scope.} New statements are numbered within this section (and within \S\ref{sec:rev8impl}); no lettered theorem is introduced, no label is ``Theorem~L'' (the framework's Parity Dominance Theorem), and none collides with L1--L7. [S38] numbers its items K-1,\dots,K-9; these are quoted for traceability only. No claim is made concerning HQCC termination, the twin prime conjecture or the Riemann hypothesis, and nothing here bears on the HQCC termination claims. CORE\_FREEZE holds: nothing in this section redefines the locked $H_s$ ($c_1=(-1,1)$, $c_2=(-4,1)$); only the label representation of Theorem~\ref{thm:F} and \S\ref{sec:family} at $s=7$ is used.

\paragraph{Notation.} Throughout, $A=\left(\begin{smallmatrix}2&7\\7&7\end{smallmatrix}\right)$ is the Gram matrix of $Q_7$ (\S\ref{sec:rev7delta}), $B(u,v)=u^{\mathsf T}Av$, and a lattice is written with its \emph{norm} $x\cdot x$ (so $Q_2=\tfrac12x\cdot x$). We write $N_1$, $N_2$ (italic) for two even lattices; they have nothing to do with the needs \textsf{(N1)}--\textsf{(N3)} of the fifth-order programme in \S\ref{sec:rev8fifth}.

%---------------------------------------------------------------------
\subsection{The even lattice $N_2$ and its discriminant form}\label{sec:rev8disc}
The form $(\Z^2,A)$ is odd ($Q_7\notin\Z$), so it is not itself a K3 lattice. Its even sublattice of index $2$ is
\[
N_2:=\{x\in\Z^2:\ x_2\equiv0\bmod2\},\qquad \mathrm{Gram}\big((1,0),(0,2)\big)=\begin{pmatrix}2&14\\14&28\end{pmatrix},\qquad\det=-140 .
\]
The other natural even version is $N_1:=(\Z^2,2A)$, with Gram matrix $\left(\begin{smallmatrix}4&14\\14&14\end{smallmatrix}\right)$, also of determinant $-140$.

\begin{proposition}[Discriminant form of $N_2$; {[S38]}, K-1]\label{prop:disc8}
\begin{enumerate}[label=(\roman*)]
\item In the basis $(1,0),(-7,2)$ of $N_2$ (base change $P=\left(\begin{smallmatrix}1&-7\\0&1\end{smallmatrix}\right)$ from $(1,0),(0,2)$), $N_2\cong\langle2\rangle\oplus\langle-70\rangle$, i.e.\ $x\cdot x=2(x_1+7y)^2-70y^2$. In particular $N_2$ is even, of signature $(1,1)$.
\item $A_{N_2}=N_2'/N_2\cong\Z/2\times\Z/70$, $|A_{N_2}|=140$, generated by $(\tfrac12,0)$ and $(0,\tfrac1{70})$ in the $\langle2\rangle\oplus\langle-70\rangle$ basis, with $q_{N_2}(x,y)=\tfrac{x^2}4-\tfrac{y^2}{140}\bmod1$.
\item The local Jordan components and their Gauss-sum signatures are:
\begin{center}\small
\begin{tabular}{llll}\toprule
$p$ & group & $q$ on the $p$-part & signature mod $8$\\\midrule
$2$ & $(\Z/2)^2$, generators $(\tfrac12,0),(0,\tfrac12)$ & $\langle\tfrac14\rangle\oplus\langle\tfrac14\rangle$, $q=(x^2+y^2)/4$ & $2$\\
$5$ & $\Z/5$ & $\langle\tfrac35\rangle\cong\langle\tfrac25\rangle$ & $4$\\
$7$ & $\Z/7$ & $\langle\tfrac27\rangle\cong\langle\tfrac17\rangle$ & $2$\\\bottomrule
\end{tabular}
\end{center}
The local signatures sum to $2+4+2\equiv0=\mathrm{sign}(N_2)\bmod8$ (oddity formula). The $5$- and $7$-parts are exactly the local forms $q_5=2x^2/5$ and $q_7=x^2/7$ of Theorem~\ref{thm:tensor7}.
\end{enumerate}
\end{proposition}
\begin{proof}
(i) $(1,0)\cdot(1,0)=2$, $(-7,2)\cdot(-7,2)=98-196+28=-70$, $(1,0)\cdot(-7,2)=-14+14=0$, and $\det P=1$. (ii) is read off from (i): $\langle2\rangle'=\tfrac12\Z$ with $q=\tfrac12\cdot2\cdot\tfrac14$, $\langle-70\rangle'=\tfrac1{70}\Z$ with $q=-\tfrac1{140}$. (iii) The $2$-part is generated by $(\tfrac12,0)$ and $(0,\tfrac{35}{70})$, with $q(0,\tfrac12)=-35^2/140\equiv\tfrac14$; the $5$-part by $(0,\tfrac15)$, $q=-14^2/140\equiv\tfrac35$; the $7$-part by $(0,\tfrac17)$, $q=-10^2/140\equiv\tfrac27$. The isomorphisms hold because $3\cdot2^{-1}\equiv4$ is a square mod $5$ and $2\equiv3^2$ mod $7$. The signatures are the standard Gauss-sum values.
\end{proof}
\catA{Elementary; exact (\texttt{closed\_exact.py} records the diagonalisation $P^{\mathsf T}\mathrm{Gram}\,P=\mathrm{diag}(2,-70)$).}

%---------------------------------------------------------------------
\subsection{The 140-label representation is the Weil representation of $N_2$}\label{sec:rev8weil}
\paragraph{Conventions.} For an even lattice $L$ we use Borcherds' convention \cite{Borcherds}: $\rho_L(T)\e_\gamma=\e(q(\gamma))\e_\gamma$ and $\rho_L(S)\e_\gamma=\e(-\mathrm{sign}(L)/8)\,|A_L|^{-1/2}\sum_\delta\e(-b(\gamma,\delta))\e_\delta$; the dual $\rho_L^*$ is the complex conjugate. Since $\mathrm{sign}(N_2)=0$, $\rho_{N_2}$ is an honest representation of $\mathrm{SL}_2(\Z)$ and $\e(\mathrm{sign}/8)=1$. On the label side, the labels are $(\alpha,\beta)$ with $\alpha\in(L'\cup(a+L'))/\Z^2$ ($70$ classes), $\beta\in\{0,\beta_1\}$, $a=(\tfrac7{10},-\tfrac15)$ and $\beta_1$ the stored label of $-a$ mod $\Z^2$; the operators $S_{\rm lab}$, $T_{\rm lab}$ are built in the convention $Y[i,j(i)]$ of Proposition~\ref{prop:audit6} (with \texttt{family\_s.Family} imported unmodified and exact rational phases, all denominators dividing $140$). Following Remark~\ref{rem:Sconv}, the label representation in standard normalisation is
\[
\rho_{\rm rev}:\quad S\mapsto S_{\rm lab}^{-1},\qquad T\mapsto T_{\rm lab}.
\]

\begin{theorem}[Weil realisation of the label representation; {[S38]}, K-2, K-4]\label{thm:weil8}
\begin{enumerate}[label=(\roman*)]
\item $\rho_{\rm rev}\cong\rho_{N_2}$ as representations of $\mathrm{SL}_2(\Z)$. No character twist is needed and no metaplectic sign arises.
\item An explicit intertwiner is the following. Index the basis of $\rho_{N_2}$ by $z\in N_2'/N_2$, written in $\Z^2$-coordinates, and let $[z]=z\bmod\Z^2$. Every class $\alpha$ has exactly two lifts $z$, $z+(0,1)$ in $N_2'/N_2$. Put
\[
X\,\e_z\;=\;u_{([z],0)}\;+\;\e\big(B(z,\beta_1)\big)\,u_{([z],\beta_1)} .
\]
Then $X\rho_{N_2}(S)=S_{\rm lab}^{-1}X$, $X\rho_{N_2}(T)=T_{\rm lab}X$, and $X$ is invertible.
\item $\rho_{\rm rev}\not\cong\rho_{N_2}^*$, $\rho_{\rm rev}\not\cong\rho_{N_1}$ and $\rho_{\rm rev}\not\cong\rho_{N_1}^*$.
\end{enumerate}
\end{theorem}
\begin{proof}
(ii) All entries of $\sqrt{140}\,X\rho_{N_2}(S)$, $\sqrt{140}\,S_{\rm lab}^{-1}X$ (note $\sqrt{140}=2\sqrt{35}$ and $\sqrt{35}\,S_{\rm lab}^{-1}$ has entries $\e(\cdot)$), $X\rho_{N_2}(T)$ and $T_{\rm lab}X$ lie in $\Z[\zeta_{280}]$. Each of the $2\cdot140^2$ entry identities is the vanishing of an explicit integer combination of powers of $\zeta_{280}$; it was checked by exact reduction modulo the cyclotomic polynomial $\Phi_{280}$, with $0$ failures for $S$ and $0$ for $T$. $X$ is block diagonal with respect to $\alpha$; the block of the two lifts $z,z'=z+(0,1)$ is $\left(\begin{smallmatrix}1&1\\\e(B(z,\beta_1))&\e(B(z',\beta_1))\end{smallmatrix}\right)$, with determinant $\e(B(z,\beta_1))\big(\e(B((0,1),\beta_1))-1\big)\neq0$ because $B((0,1),\beta_1)=7(-\tfrac7{10})+7\cdot\tfrac15\equiv\tfrac12\bmod1$. (i) follows from (ii); $X$ intertwines the untwisted $\rho_{N_2}$, and $\mathrm{sign}(N_2)=0$.
(iii) $T_{\rm lab}$ is monomial; its label permutation has $70$ fixed points and $35$ two-cycles, and its eigenvalue multiset ($36$ distinct values) equals that of $\rho_{N_2}(T)$ but differs (exactly, as multisets of elements of $\Q/\Z$) from those of $\rho_{N_2}^*(T)$, $\rho_{N_1}(T)$ and $\rho_{N_1}^*(T)$. Conjugate representations have conjugate $T$-spectra.
\end{proof}
\catA{Exact computation in $\Z[\zeta_{280}]$ (\texttt{closed\_exact.py} $\to$ \texttt{closed\_exact.json}: \texttt{exact\_S\_failures}$=0$, \texttt{exact\_T\_failures}$=0$, $B(v,\beta_1)\equiv\tfrac12$, invertible); exact $T$-spectrum comparison (\texttt{lattices.py}). A second exact solution, with offsets $\{0,\tfrac45\}$, is in \texttt{explicit.py}/\texttt{X\_explicit.pkl}; it is another element of the Hom space of Proposition~\ref{prop:hom8}. Further $T$-spectrum tests: $2\langle5,-7\rangle$, $2\langle-5,7\rangle$, $2\left(\begin{smallmatrix}3&1\\1&-12\end{smallmatrix}\right)$ and its negative do not match; $2\langle1,-35\rangle$ and $2\left(\begin{smallmatrix}1&1\\1&-34\end{smallmatrix}\right)$ match only because both are $\cong N_2$, and their negatives match the dual.}

\begin{proposition}[Relations and Hom spaces modulo $\ell$; {[S38]}, K-3]\label{prop:hom8}
Let $\ell=840841\equiv1\bmod280$. Computing modulo $\ell$ (with a fixed primitive $280$-th root of unity in $\mathbb F_\ell$):
\begin{enumerate}[label=(\roman*)]
\item $(S_{\rm lab}^{-1},T_{\rm lab})$ satisfies $S^4=I$ and $(ST)^3=S^2$, as $\rho_{N_2}$ does; $(S_{\rm lab},T_{\rm lab})$ satisfies $S^4=I$ but \emph{fails} $(ST)^3=S^2$, so it is not a representation of $\mathrm{SL}_2(\Z)$ in this normalisation (confirming Remark~\ref{rem:Sconv}).
\item $\dim_{\mathbb F_\ell}\mathrm{Hom}(\rho_{N_2},(S_{\rm lab}^{-1},T_{\rm lab}))=8$, and a random element has rank $140$.
\item $\dim_{\mathbb F_\ell}\mathrm{Hom}(\rho_{N_2},(S_{\rm lab},T_{\rm lab}))=4$, and a random element has rank $72$: the two normalisations agree only on the $72$-dimensional subspace $W_P$.
\end{enumerate}
\end{proposition}
\catA{Exact computation modulo $\ell$ (\texttt{intertwine.py} $\to$ \texttt{intertwine.json}, \texttt{X\_modl.pkl}). The isomorphism itself is Theorem~\ref{thm:weil8}, which supersedes (ii) for that purpose. By semicontinuity the characteristic-$0$ Hom dimensions are at most $8$ and $4$. \emph{Category~B (reading):} the value $8$ is consistent with eight pairwise non-isomorphic irreducible sectors $W_5^\pm\otimes W_7^\pm\otimes\{A_3,B_1\}$ (\S\ref{sec:family}); that the characteristic-$0$ dimension equals $8$ is not separately proved here.}

\begin{remark}[Dual, twist, metaplectic sign]\label{rem:twist8}
$\rho_{\rm rev}\not\cong\rho_{N_2}^*$ (Theorem~\ref{thm:weil8}(iii)); no twist by a character of $\mathrm{SL}_2(\Z)$ is needed; $\mathrm{sign}(N_2)=0$, so there is no metaplectic sign. The Weil representation depends only on the discriminant form, so any even lattice with discriminant form $\cong q_{N_2}$ and signature $\equiv0\bmod8$ (e.g.\ $N_2\oplus U$, $N_2\oplus E_8$, and the lattice $T$ of Proposition~\ref{prop:T8}, up to dualisation) gives the same representation. The earlier alternative ``possibly $\rho_L$ or $\rho_L^*$'' ([S39], \S4(b)2 and (d)) is resolved: it is $\rho_L$ for $L=N_2$.
\end{remark}

%---------------------------------------------------------------------
\subsection{The transcendental lattice}\label{sec:rev8T}
Let $\Lambda_{K3}=U^3\oplus E_8(-1)^2$ (even unimodular, signature $(3,19)$), and let $U\oplus U\subset\Lambda_{K3}$ have bases $e_1,f_1,e_2,f_2$ with $e_i\cdot f_i=1$ and all other products $0$.
\begin{proposition}[Explicit primitive embedding and $T$; {[S38]}, K-5]\label{prop:T8}
Put $n_1=e_1+f_1$ (norm $2$), $n_2=e_2-35f_2$ (norm $-70$), $t_1=e_1-f_1$ (norm $-2$), $t_2=e_2+35f_2$ (norm $70$).
\begin{enumerate}[label=(\roman*)]
\item $\mathrm{span}(n_1,n_2)\cong N_2$ and is primitive in $U\oplus U$ (the $2\times2$ minors of $(n_1\ n_2)$ have gcd $1$).
\item $t_1,t_2$ are orthogonal to $n_1,n_2$ and span the full orthogonal complement of $\mathrm{span}(n_1,n_2)$ in $U\oplus U$ (index $1$; their Gram determinant is $-140$).
\item Hence the orthogonal complement of $N_2$ in $\Lambda_{K3}$ is
\[
T:=N_2^\perp=\langle-2\rangle\oplus\langle70\rangle\oplus U\oplus E_8(-1)^2,\qquad\mathrm{sign}(T)=(2,18),\qquad A_T\cong\Z/2\times\Z/70,
\]
with $q_T(x,y)=-\tfrac{x^2}4+\tfrac{y^2}{140}=-q_{N_2}(x,y)$ under the identity identification of generators. The gluing is the one induced by the unimodular overlattice $U\oplus U$ (consistent with Nikulin \cite{Nikulin}, Prop.~1.6.1).
\end{enumerate}
\end{proposition}
\begin{proof}
Direct computation: $n_1\cdot n_1=2$, $n_2\cdot n_2=-70$, $n_1\cdot n_2=0$, $t_1\cdot t_1=-2$, $t_2\cdot t_2=70$, and $n_i\cdot t_j=0$. The complement of a rank-$2$ primitive sublattice of the unimodular $U\oplus U$ has rank $2$ and $|\det|=140$; $\mathrm{span}(t_1,t_2)$ lies in it with $|\det|=140$, so the index is $1$. The other summands $U\oplus E_8(-1)^2$ are orthogonal to $U\oplus U$. The discriminant form of $\langle-2\rangle\oplus\langle70\rangle$ is read off as in Proposition~\ref{prop:disc8}.
\end{proof}
\catA{Elementary; exact (\texttt{closed\_exact.json}: norms $[2,-70,-2,70]$, orthogonality $[0,0,0,0]$, gcd of minors $1$, complement index $1$). Nikulin's existence theorem (Cor.~1.12.3) is \emph{not} needed, since the embedding is explicit.}

\begin{corollary}[The residual law in K3 terms; {[S38]}, K-6]\label{cor:T8}
$\rho_{\rm rev}\;\cong\;\rho_{N_2}\;\cong\;\rho_T^*=\overline{\rho_T}$ (Borcherds' convention; honest $\mathrm{SL}_2(\Z)$-representations).
\end{corollary}
\begin{proof}
$\mathrm{sign}(T)=-16\equiv0\bmod8$ and $q_T=-q_{N_2}$ (hence $b_T=-b_{N_2}$) under the identification of Proposition~\ref{prop:T8}, so in the formulas for $\rho_T(S)$ and $\rho_T(T)$ every phase is the conjugate of the corresponding phase for $\rho_{N_2}$. Combine with Theorem~\ref{thm:weil8}.
\end{proof}
\catA{Purely lattice-theoretic; no cited geometric theorem is used. In words: the $140$-label residual law transforms under the dual Weil representation of a signature-$(2,18)$ lattice, which (Theorem~\ref{thm:K3exist8}) is the transcendental lattice of an $N_2$-polarised K3 surface. This is a statement about a shared \emph{representation}; it is not a statement that any residual series is a K3 generating function (see Remark~\ref{rem:open8}).}

%---------------------------------------------------------------------
\subsection{Geometric realisation}\label{sec:rev8geom}
\begin{theorem}[$N_2$-polarised K3 surfaces; {[S38]}, K-7]\label{thm:K3exist8}
\begin{enumerate}[label=(\roman*)]
\item The primitive embedding $N_2\hookrightarrow\Lambda_{K3}$ is unique up to $O(\Lambda_{K3})$; hence $T$ is unique up to isometry.
\item There is a projective K3 surface $X$ with $\mathrm{NS}(X)\cong N_2$ and $T_X\cong T$.
\item $N_2$ represents neither $-2$ nor $0$ (non-trivially). Hence such an $X$ contains no smooth rational curves and admits no elliptic fibration, and its ample cone is the whole component of the positive cone containing an ample class.
\end{enumerate}
\end{theorem}
\begin{proof}[Proof sketch]
(i) Nikulin \cite{Nikulin}, Thm.~1.14.4: $\Lambda_{K3}$ is even unimodular of signature $(3,19)$, $N_2$ is even of signature $(1,1)$, $3>1$, $19>1$, and $\mathrm{rk}\,\Lambda_{K3}-\mathrm{rk}\,N_2=20\ge2+l(A_{N_2})=4$. (ii) Morrison \cite{Morrison}, Cor.~2.9(i): for $\rho\le10$ every even lattice of signature $(1,\rho-1)$ is the N\'eron--Severi lattice of some algebraic K3 surface (equivalently, surjectivity of the period map); with (i), $T_X\cong T$. (iii) $x\cdot x=2\big((x_1+7y)^2-35y^2\big)$; $-2$ would need $u^2-35y^2=-1$, impossible since $-1$ is not a square mod $7$; $0$ would need $35$ to be a square. Without $(-2)$-classes the ample cone is a full component of the positive cone, and an elliptic fibration requires a non-zero isotropic class in $\mathrm{NS}$ (standard, \cite{Huybrechts}).
\end{proof}
\catA{\textbf{Category~A (cited)}: conditional on Nikulin 1.14.4 and Morrison Cor.~2.9(i) (and standard facts on ample cones and elliptic fibrations, \cite{Huybrechts}). The arithmetic in (iii) is exact (\texttt{checks.json}, IC-5: PARI \texttt{qfbsolve} finds no solutions). \emph{Citation status:} the theorem numbers were checked against secondary sources (arXiv:1312.4421, arXiv:2303.03294, and a Springer paper on coupling pairs), not against the original articles.}

%---------------------------------------------------------------------
\subsection{Automorphisms}\label{sec:rev8aut}
Let $\gamma=\left(\begin{smallmatrix}-1&-7\\2&13\end{smallmatrix}\right)$, the Pell automorph of Theorem~\ref{thm:B}: $\gamma^{\mathsf T}A\gamma=A$, $\det\gamma=1$, $\mathrm{tr}\,\gamma=12$, eigenvalues $6\pm\sqrt{35}$. It preserves $N_2$ (the second coordinate of $\gamma x$ is $2x_1+13x_2$). Let $r=\mathrm{diag}(1,-1)$ in the $\langle2\rangle\oplus\langle-70\rangle$ basis; there $\gamma=\left(\begin{smallmatrix}6&35\\1&6\end{smallmatrix}\right)$ (multiplication by $\varepsilon=6+\sqrt{35}$).

\begin{proposition}[Action on $A_{N_2}$ and $O(N_2)$; {[S38]}, K-8]\label{prop:gamma8}
\begin{enumerate}[label=(\roman*)]
\item The induced map $\bar\gamma$ on $A_{N_2}$ acts on the local parts as follows: on the $2$-part $(\Z/2)^2$ as a non-trivial involution (not the identity, which on this part equals $\pm1$); on the $5$-part as $+1$; on the $7$-part as $-1$. Hence $\bar\gamma\neq\pm1$ and $\bar\gamma^2=\mathrm{id}$. This matches the local form $Z_{(s)}\otimes I\otimes Z_2$ on the $7$/$5$/$2$ parts in \S\ref{sec:rev6}.
\item $O(N_2)=\{\pm\gamma^k,\ \pm r\gamma^k:\ k\in\Z\}$, and $r\gamma r=\gamma^{-1}$.
\item On $A_{N_2}$: $\gamma^{2j}$ acts as $+1$ and $r\gamma^{2j}$ as $-1$, both preserving the positive cone; odd powers $\pm\gamma^{2j+1}$, $\pm r\gamma^{2j+1}$ act as neither $+1$ nor $-1$; $-\gamma^{k}$ and $-r\gamma^k$ do not preserve the positive cone.
\end{enumerate}
\end{proposition}
\begin{proof}
(i), (iii): exact computation on the $140$ classes (\texttt{closed\_exact.py}, \texttt{aut.py}). (ii): in the diagonal basis $x\cdot x=2(u^2-35y^2)$. Proper automorphs of $u^2-35y^2$ correspond to solutions of $t^2-140w^2=4$, i.e.\ to norm-$1$ units of $\Z[\sqrt{35}]$, which are $\pm\varepsilon^k$ with $\varepsilon=6+\sqrt{35}$ fundamental; there is no unit of norm $-1$ ($u^2-35y^2=-1$ is impossible mod $7$). Improper automorphs are $r$ times proper ones, and $r\gamma r=\left(\begin{smallmatrix}6&-35\\-1&6\end{smallmatrix}\right)=\gamma^{-1}$.
\end{proof}
\catA{Exact (\texttt{closed\_exact.json}: \texttt{gamma\_on\_A\_N2\_local} $=\{2:\text{neither},\,5:+1,\,7:-1\}$, \texttt{gamma\^{}2\_is\_+1}; \texttt{aut.json}). \emph{Correction to [S39], IC-6:} [S39] also tabulated $\gamma$ on $A_{N_1}$; that computation is correct as arithmetic, but it is \emph{not} used, because the label representation is $\rho_{N_2}$ and not $\rho_{N_1}$ (Theorem~\ref{thm:weil8}(iii)).}

\begin{theorem}[Automorphisms of $N_2$-polarised K3 surfaces; {[S38]}, K-8]\label{thm:aut8}
Let $X$ be as in Theorem~\ref{thm:K3exist8}, with $\mathrm{NS}(X)=N_2$ glued to $T_X=T$ as in Proposition~\ref{prop:T8}.
\begin{enumerate}[label=(\roman*)]
\item $\gamma^2\oplus\mathrm{id}_{T}$ extends to a Hodge isometry of $H^2(X,\Z)$ preserving the ample cone, and is induced by a unique $f\in\mathrm{Aut}(X)$. $f$ has infinite order; $f^*$ has spectral radius $(6+\sqrt{35})^2=71+12\sqrt{35}$ on $H^2(X,\Z)$ (a quadratic unit, not a Salem number of higher degree), so its topological entropy is
\[
h(f)=2\log(6+\sqrt{35})=4\sqrt{35}\,I\approx4.9558,
\]
where $I=\log(6+\sqrt{35})/(2\sqrt{35})$ is the constant of Theorem~\ref{thm:C}.
\item $r\oplus(-\mathrm{id}_T)$ is induced by a non-symplectic involution $\iota\in\mathrm{Aut}(X)$.
\item If $X$ is very general in the $N_2$-polarised family, so that the Hodge isometries of $T_X$ are only $\pm1$, then $\mathrm{Aut}(X)=\langle f,\iota\rangle\cong\Z\rtimes\Z/2$ (the infinite dihedral group $D_\infty$), and $\gamma$ itself does not lift to $\mathrm{Aut}(X)$.
\end{enumerate}
\end{theorem}
\begin{proof}[Proof sketch]
By Nikulin's gluing criterion \cite{Nikulin}, Cor.~1.5.2, an isometry $g_{\rm NS}\oplus g_T$ of $N_2\oplus T$ extends to $H^2(X,\Z)$ iff $\bar g_{\rm NS}$ and $\bar g_T$ agree under the glue isomorphism $A_{N_2}\cong A_T$. By Proposition~\ref{prop:gamma8}, $\bar\gamma^2=+1$ and $\bar r=-1$, compatible with $\mathrm{id}_T$ and $-\mathrm{id}_T$, which are Hodge isometries. Both preserve the positive cone, hence (Theorem~\ref{thm:K3exist8}(iii)) the ample cone. By the strong Torelli theorem (Piatetski-Shapiro--Shafarevich, Burns--Rapoport; e.g.\ \cite{Huybrechts}, Thm.~7.5.3) they are induced by unique automorphisms $f$, $\iota$; $\iota$ acts by $-1$ on $H^{2,0}$, so it is non-symplectic. The eigenvalues of $f^*$ are $(6\pm\sqrt{35})^2$ on $\mathrm{NS}\otimes\R$ and $1$ on $T$, and the entropy is the logarithm of the spectral radius (Gromov--Yomdin \cite{GY}). For (iii): every automorphism acts on $T_X$ by $\pm1$, so its restriction to $N_2$ preserves the positive cone and acts on $A_{N_2}$ by $\pm1$; by Proposition~\ref{prop:gamma8}(iii) these restrictions are exactly $\gamma^{2j}$ and $r\gamma^{2j}$, which form $\langle\gamma^2\rangle\rtimes\langle r\rangle\cong D_\infty$; the action on $H^2$ is injective by Torelli. Since $\bar\gamma\neq\pm1$, $\gamma$ does not lift.
\end{proof}
\catA{\textbf{Category~A (cited)}: conditional on the strong Torelli theorem, Nikulin's gluing (Cor.~1.5.2), the Gromov--Yomdin entropy formula, and the standard fact that a very general member has only $\pm1$ as Hodge isometries of $T_X$. The lattice inputs (Proposition~\ref{prop:gamma8}) are exact. The numerical value $\approx4.9558$ is $2\log(11.91608\ldots)$; [S38] prints ``$\approx4.93$'', which is a typo (corrected here). The identity $2\log(6+\sqrt{35})=4\sqrt{35}\,I$ is a definitional identity; \emph{Category~B (reading):} that the regulator driving the Baker obstruction is ``the same'' as this entropy is an observation about definitions, and no theorem linking $C$ to $X$ is claimed.}

\begin{remark}[Category~B and open items for \S\ref{sec:rev8k3}]\label{rem:open8}
None of the following is a result of this paper.
\begin{enumerate}[label=(\alph*)]
\item \emph{Noether--Lefschetz / Heegner series} ([S38], K-9; Category~B). The Noether--Lefschetz/Heegner generating series of the $N_2$-polarised family are vector-valued modular forms of weight $1+18/2=10$ for $\rho_T$ or $\rho_T^*$ (depending on convention; Borcherds, Kudla--Millson, Maulik--Pandharipande, cited by name only). The residual objects of Theorems~\ref{thm:F}--\ref{thm:G} are weight-$0$ Maass-type objects with $\lambda=\tfrac14$. A shared \emph{representation} (Corollary~\ref{cor:T8}) is not a shared \emph{form}; no theta-lift or Borcherds-product link has been checked.
\item \emph{Special members} (open). Whether some special $N_2$-polarised $X$ has a Hodge isometry $\psi$ of $T_X$ with $\bar\psi=\bar\gamma$, so that $\gamma$ itself lifts. This needs the Hodge isometries of $T_X$ to form a cyclic group of order $m>2$ with $\varphi(m)\mid20$, plus compatibility of $\bar\psi$. Not checked.
\item \emph{Interpretation} (Category~B). Any physical or ``unification'' reading of $X$, $f$ or $\iota$.
\item \emph{Citation hygiene} (open). The theorem numbers for Nikulin (1979/80) and Morrison (1984) should be confirmed against the original articles. The ``automorphs of binary forms $\leftrightarrow$ K3 automorphisms'' literature (Galluzzi--Lombardo--Peters, cited in [S39] from memory) is Category~B as a citation.
\item \emph{Scope} (correction to [S39], \S4 and IC-6). [S39] treated both $N_1$ and $N_2$ as K3 candidates and also stated K3 consequences for the lattices $M_t$ of Proposition~K1. In this paper every K3 statement is made for $N_2$ only, because the $140$-label representation is $\rho_{N_2}$ and not $\rho_{N_1}$ (Theorem~\ref{thm:weil8}(iii)). The exact lattice facts of IC-5 for $N_1$ and $M_t$ (anisotropic, no $-2$ vectors) remain true but are not used.
\end{enumerate}
\end{remark}

%=====================================================================
\section{Implications}\label{sec:rev8impl}
%=====================================================================
This section (revision~8) records what the Bridge-C results of this paper imply for four parts of the author's wider programme, following [S39] (\path{2026-10-19-Implications-Unification-RTA-multik-K3.md}, with checks IC-1--IC-7 in \path{2026-10-19-Implications-checks/}), as corrected by [S38], \S6. The four terms are defined in the author's corpus (\path{docs/S2-11DM2ET-X-Axiomatic-Book/}, read-only): ``fifth-order unification'' (the exact phrase ``Order of Unification'' does not occur), the Residual Ternary Algebra manuscript (Ramanujan Journal v1.0, submitted 2026-08-26), ``pure-even multi-$k$'' families (Bring--Jerrard/$A_5$), and K3 material that the corpus itself labels Category~B. [S39] was written against revision~6; the theorem names it uses (A--K, G$_s$) are unchanged in this revision.
Category convention as in \S\ref{sec:cat}, and the two categories are kept strictly separate: facts are labelled Category~A only where they are proved or exact; every reading, analogy or identification is Category~B. No Category~B item is imported as a theorem (in particular Theorem~\ref{thm:K} and the PSLQ-only items remain Category~B). No claim is made concerning HQCC termination, the twin prime conjecture or the Riemann hypothesis; this section is separate from the HQCC termination claims and uses no lepton, Orch-OR or neuroscience material. CORE\_FREEZE holds: the locked $H_s$ with $c_1=(-1,1)$, $c_2=(-4,1)$ is not redefined; $H_s^{\rm axes}$ (\S\ref{sec:rev7axes}) is a parallel object.

%---------------------------------------------------------------------
\subsection{Residual Ternary Algebra}\label{sec:rev8rta}
The RTA manuscript (\S\S9--12) states that ``the residual law for $C$ --- the exact weight-$\tfrac12$ slash prefactor identity and quantitative continuous Fricke verification against residual samples --- remains open''.

\begin{proposition}[No scalar residual law; restatement of Theorem~\ref{thm:E}]\label{prop:noscalar8}
\begin{enumerate}[label=(\roman*)]
\item There is no identity $C=F+\mathcal E$ with $F$ an algebraic combination of holomorphic modular forms, signature-$(1,1)$ mock (Zwegers) holomorphic parts, unary, partial or false theta series, and $\mathcal E(e^{-t})=o(1/t)$. The reason is the leading term $C(e^{-t})=I/t+a_Ct^{-1/2}+\cdots$ with $I=\log(6+\sqrt{35})/(2\sqrt{35})$ (Theorem~\ref{thm:C}) and $\log(6+\sqrt{35})\notin\Qbar+\pi\Qbar$ (Baker).
\item There is no quotient law $C\cdot g=F$ with holomorphic modular $g\not\equiv0$, $F$.
\item There is no scalar $W_7$ law: for no weight $k$ and no constant $\lambda$ is $C(\tau)-\lambda(7\tau)^{-k}C(-1/(7\tau))$ radially bounded at all rationals.
\end{enumerate}
In particular the ``exact weight-$\tfrac12$ slash prefactor identity'' of RTA \S12 does not exist; RTA's own finding ``residual constant $\times$ residual automorphy therefore fails'' is confirmed and proved.
\end{proposition}
\catA{(i) proved \emph{conditional on} Baker's theorem \cite{Baker} and the $q$-expansion principle (both cited). (ii) proved (for $C$ using the closed form of $a_C\neq0$, modulo F1). (iii) is an \emph{exact finite computation}: at $178$ of the $270$ test rationals (denominator $<30$), including the Master sample $x=5/7$, exactly one of $G(x)$, $G(-1/(7x))$ vanishes (exact by Theorem~\ref{thm:D}). See \S\ref{sec:baker}.}

\begin{proposition}[The law that exists has weight $1$ and is vector-valued]\label{prop:weight8}
\begin{enumerate}[label=(\roman*)]
\item Where both Gauss means are nonzero, $|G|$ obeys an exact weight-$1$ relation with $|\lambda|=7^{-1/2}$ (\S\ref{sec:baker}); the leading $1/t$ radial profiles obey the $S/T$ law with factor $1/|x|$ (Theorem~\ref{thm:profile}); and the vector cocycle has automorphy factor $|cx+d|$ (Theorem~\ref{thm:G}).
\item The replacement for the scalar residual law is the vector-valued weight-$1$, $140$-component law: the $140$-label Siegel lift to $\lambda=\tfrac14$ Maass forms (Theorem~\ref{thm:F}), the Lewis--Zagier period formula (Theorem~\ref{thm:G}), the $42$-dimensional function span $V=V_{\rm Eis}(36)\oplus V_{\rm cusp}(6)$ (Theorem~\ref{thm:H}), the Herglotz closed form (Theorems~\ref{thm:I}, \ref{thm:J}, given (I1)--(I3)), and, since this revision, the identification $\rho_{\rm rev}\cong\rho_{N_2}$ (Theorem~\ref{thm:weil8}). It is a law for the vector of cone sums, not for scalar $C$.
\item RTA \S10's characteristic $\chi_{\rm nat}(s)=\big(\tfrac s{2(s-2)},-\tfrac1{s-2}\big)$ is the characteristic $a=A^{-1}(0,\tfrac s2)^{\mathsf T}$ that completes the square: $Q_+(v)=Q_2(v+a)+c_0$ with $c_0=\tfrac s{4(s-2)}$.
\end{enumerate}
\end{proposition}
\catA{(i), (ii): with the status of the cited theorems as stated where they appear (in particular Theorems~\ref{thm:I}, \ref{thm:J} are given (I1)--(I3); Theorem~\ref{thm:K} is \emph{not} part of the law and stays Category~B). (iii): $A^{-1}=\frac{-1}{s(s-2)}\left(\begin{smallmatrix}s&-s\\-s&2\end{smallmatrix}\right)$ gives $a=\big(\tfrac s{2(s-2)},-\tfrac1{s-2}\big)$, $B(v,a)=\tfrac s2k$, and $-Q_2(a)=-\tfrac12a\cdot(0,\tfrac s2)=\tfrac s{4(s-2)}$. Consequently the weight $\tfrac12$ of RTA \S\S11--12 is the wrong weight for the law that exists.}

\begin{remark}[The samples are not radial values]\label{rem:samples8}
The four published ``residual samples'' ($19.80$, $1.82$, $9.24$, $20.90$ at $x=1/7,3/7,5/7,1$) are not radial values of $C$: $|C(0.995\,\e(x))|\approx19.5,18.5,19.9,54.0$. They are replaced by the exact radial targets $\lim_{t\to0}t\,C(\e(x)e^{-t})=I\cdot G(x)$ at every rational $x$ (Theorem~\ref{thm:D}). So ``quantitative continuous Fricke verification against residual samples'' is not a meaningful target, and this paper does not claim it.
\end{remark}
\catB{The four moduli at $0.995\,\e(x)$ are floating-point evaluations. The replacement targets $I\cdot G(x)$ are Category~A (Theorems~\ref{thm:C}, \ref{thm:D}).}

\begin{proposition}[Two factual corrections to the RTA text; IC-1, IC-2]\label{prop:ic12}
With $s=7$, $c_1=(-1,1)$, $c_2=(-4,1)$ and $C=P_s-H_s$ as in \S\ref{sec:main}:
\begin{enumerate}[label=(\roman*)]
\item (IC-1.) The coefficients of $C$ lie in $\{0,1,\tfrac32\}$ \emph{only through} $q^{101}$. The first coefficient outside this set is $[q^{102}]C=2$ (coming from $\Cthird$); the first eight exponents with a coefficient outside the set are $102,121,144,170,183,196,253,289$, with value $2$ except $[q^{144}]C=3$. Through $q^{419}$ the value set is $\{0,\tfrac12,1,\tfrac32,2,3\}$, the values $\tfrac12$ coming from $\Cmixed=-\tfrac12q^{294}+\tfrac12q^{301}-q^{330}+q^{337}-q^{368}+q^{375}-\cdots$ (reproducing Theorem~\ref{thm:A}(iv)). RTA \S9's ``through degree $100$ the coefficients of $C$ take values only in $\{0,1,1.5\}$'' is literally true but is a finite-range observation, not a structural property; this agrees with item~4 of \S\ref{sec:recon}.
\item (IC-2.) RTA uses the name $P_s$ for two different series: \S7 sets $P_s:=(-q^s;q^s)_\infty\,g_s$, while \S9 uses $P_s$ for the orthant series $\sum_{n,k\ge0}q^{Q_s(n,k)}$. At $s=7$ they first differ at $q^{14}$ (coefficients $1$ and $0$ respectively; further differences at $q^{21},q^{22},q^{28},\dots$), because of the extra factor $1/(q^s;q^s)_k$. This paper uses $P_s$ only in the \S9 (orthant) sense; a revision of the RTA text should rename one of the two.
\end{enumerate}
\end{proposition}
\catA{Exact rational arithmetic (\texttt{cvals.py} $\to$ \texttt{cvals.json}, $N=420$; \texttt{checks.py} $\to$ \texttt{checks.json}, key \texttt{RTA\_two\_Ps\_first\_diffs\_s7}).}

\begin{remark}[RTA \S11: ``holomorphic projection recovers $C$'']\label{rem:holproj8}
Any weight-$\tfrac12$ harmonic Maass form has radial growth $O(t^{-1/2})$, so no weight-$\tfrac12$ object can carry the $I/t$ term of $C$; whatever the formal $\Xi$ and its Eichler integral are, they can at most describe pieces of $C$, not its residual law. The growth argument is Category~A, conditional on the standard (cited) $t^{-1/2}$ radial bound. \emph{Category~B:} the agreement of RTA's ``disc-$35$'' $\Xi$ with the field $\Q(\sqrt{35})$ of Theorem~\ref{thm:B} is an identification only, since $\Xi$ is defined only formally there.
\end{remark}

\noindent\textbf{Unaffected and open.} RTA \S\S2--6, \S8 and \S13 (Zwegers completion of $H_s$, admissible for integer $s\ge5$) are unaffected. Open: uniqueness or normalisation of the vector completion; quantum modularity of scalar $C$ and of $g_s$ (not claimed: we do not claim that $C$, $\Cthird$ or $P_s$ is modular, mock modular or quantum modular as a scalar $q$-series); sharpness of the Proposition~K1 bounds at $m=17,31,\dots$; $t^{19/2}$ and beyond. As in [S39], the submitted RTA v1.0 is left as is; this section is internal guidance for a possible requested revision only.

%---------------------------------------------------------------------
\subsection{Fifth-order unification}\label{sec:rev8fifth}
The corpus requires three ingredients (\path{Executive_Summary_2026-09-17_Needed_Ingredients.md}; glossary): \textsf{(N1)} a coefficient-emitting map from $E_8$/PW data to the series (an $E_8$ functor is \emph{not locked}); \textsf{(N2)} a unique completed period for $C$; \textsf{(N3)} a tail for $C$. (These \textsf{(N1)}--\textsf{(N3)} are unrelated to the lattices $N_1$, $N_2$ of \S\ref{sec:rev8k3}.)
\begin{center}\small
\begin{longtable}{p{2.2cm}p{8.3cm}p{4.2cm}}\toprule
need & status after this paper & category\\\midrule\endhead
\textsf{(N1)} map/functor & No bearing: no result here touches $E_8$, PW$(7,11;3^5)$ or any functor. Unchanged (``No functor''; ``$E_8$ functor: not locked''). & ---\\
\textsf{(N2)} completed period & Scalar version excluded (Theorem~\ref{thm:E}); supplied only as the $140$-component vector cocycle (Theorems~\ref{thm:F}, \ref{thm:G}; Herglotz form Theorems~\ref{thm:I}, \ref{thm:J} given (I1)--(I3)). \emph{Uniqueness} is not addressed by any theorem: open. & A (as labelled at each theorem); uniqueness open\\
\textsf{(N3)} tail & Exact closed forms at $x=0$ through $t^{17/2}$ (\S\ref{sec:rev7}), at $x=\tfrac12,\tfrac13$ through $t^{13/2}$, the period-$72$ theorem at $x=\tfrac13$, and the Gauss evaluation and CFP at general $x=a/N$. Not given at all rationals in closed form; $t^{19/2}$ onward not computed. & A (as labelled in \S\S\ref{sec:rev6}--\ref{sec:rev7})\\\bottomrule
\end{longtable}
\end{center}

\begin{proposition}[The role of $s=7$ comes from the frozen $c_2$]\label{prop:s7role8}
$7s-32=-2Q(c_2)$ for the locked edge $c_2=(-4,1)$ (Proposition~\ref{prop:delta7}); only the factor $s-2$ of $\Delta$ is intrinsic given $c_1$ (Remark~\ref{rem:intrinsic7}). The Bridge-C mathematics is otherwise uniform in $s$: the tensor theorem holds for all odd $s$ (Theorem~\ref{thm:tensorodd}), the radial closed forms are rational functions of $s$, and Theorem~G$_s$ holds for all real $s>0$. The values that are structurally special in this work are $s=5$ (where the radial coefficients of \S\ref{sec:rev7} vanish) and $s=8$ (frozen cone $=$ axes cone, $\Delta=144=12^2$), not $s=7$ or $s=11$.
\end{proposition}
\catA{The facts are Category~A as proved in \S\S\ref{sec:hsct}, \ref{sec:rev7}. \emph{Category~B (reading):} the closed-form evidence does not single out the packaging values $s\in\{7,8,11\}$; the discrete bridge $E_8\to$ PW $\to\{7,8,11\}\to g_s$ is a packaging statement and is \textbf{not refuted}, but the Bridge-C results give it \textbf{no independent support}.}

\noindent\textbf{The $f_0$ contact.} $g_1=f_0$ by definition (Hickerson--Mortenson \cite{HM}). At $s=1$ the lattice $\left(\begin{smallmatrix}2&1\\1&1\end{smallmatrix}\right)$ has $\det=-s(s-2)=1>0$ and is positive definite, so the indefinite part of the construction ($H_s$, $C$, $\Q(\sqrt{s(s-2)})$) does not specialise to the classical fifth-order Hecke-type structure over $\Q(\sqrt5)$; $H_s$ is Zwegers-admissible only for integer $s\ge5$. The fifth-order link is carried only by $g_s$ at $s=1$, where Theorem~G$_s$ gives a proved radial expansion uniform in $s$ (Category~A). Quantum modularity of $g_s$ remains a conjecture of the corpus and is not claimed.
\catA{Elementary for the definiteness statement; Theorem~G$_s$ as labelled in Appendix~\ref{app:gs}.}
\noindent\textbf{Wording that needs correction in the corpus.} (a) Any statement that the fifth-order bridge needs a \emph{scalar} weight-$\tfrac12$ slash identity for $C$: that route is closed (Proposition~\ref{prop:noscalar8}(iii)). (b) Any implication that the $s=7$ closed forms are intrinsic to $s=7$ (Proposition~\ref{prop:s7role8}).
\noindent\textbf{Open.} \textsf{(N1)} or its retirement; uniqueness of the vector completion; closed-form tails at general rationals and from $t^{19/2}$ on; quantum modularity of $g_s$; whether an $s$-canonical cone should replace the frozen $c_2$ (an author's definitional choice; CORE\_FREEZE keeps the frozen cone).

%---------------------------------------------------------------------
\subsection{Even multi-$k$}\label{sec:rev8multik}
In the corpus, ``pure-even multi-$k$'' families are Bring--Jerrard quintics $x^5+\alpha(m)x+\beta(m)$ with $\alpha(m)=256m^2-3125k^4/256$, $\beta(m)=k\,\alpha(m)$, whose discriminant is a square (Galois group in $A_5$); the locked success criterion asks for a continuous monodromy representation of multi-$k$ $A_5$ type that forces Necessity.
\begin{enumerate}[label=(\alph*),leftmargin=2em]
\item \emph{Direct bearing: none.} No result here concerns Bring--Jerrard quintics, Hurwitz spaces or Galois groups of quintic fibres.
\item \emph{Name collisions, not identifications.} ``Even'' here means even $s$ or the $(+1)$-parts of Weil representations; in multi-$k$ it means square discriminant. The $k$ of Proposition~K1 ($m=2k^2-1$, $s=2m^2$) is a Pell index; in multi-$k$ it is a slice parameter. If ``even multi-$k$'' were read as the even-$s$, all-$k$ family of Proposition~K1 (a reading, not a corpus definition), its status is that of \S\ref{sec:rev7K1}: proved for all $k\ge2$, dimension bounds proved as upper bounds, sharpness open.
\item \emph{Formal $A_5$ contact.} $\mathrm{PSL}_2(\mathbb F_5)\cong A_5$ (standard). At $s=7$ the label representation has the tensor factor $\rho_{2x^2/5}$ (Theorem~\ref{thm:tensor7}); $V_{\rm Eis}=A_3\otimes W_5^+\otimes W_7^+$ contains $W_5^+$ (dimension $3$, factoring projectively through $A_5$) and $V_{\rm cusp}=B_1\otimes W_5^-\otimes W_7^-$ contains $W_5^-$ (dimension $2$, a faithful representation of $\mathrm{SL}_2(\mathbb F_5)=2\cdot A_5$).
\end{enumerate}
\catA{(c) as standard theory plus Theorem~\ref{thm:tensor7} and Proposition~\ref{prop:matched}. \emph{IC-4} (exact; \texttt{checks.py}, sympy and PARI): the corpus' Theorem~3.1 discriminant identity holds symbolically, the four seed discriminants are squares, and PARI \texttt{polgalois} returns $A_5$, $A_5$, $A_5$, $D_5$ for $x^5-55x+88$, $x^5-100x+400$, $x^5+20x+16$, $x^5+20x-32$. This verifies catalogue base facts only, not density or Necessity.}
\catB{\emph{IC-3} (floating-point group closure, \texttt{weil5.py} $\to$ \texttt{weil5.json}): the $S,T$ images on $W_5^+$ and $W_5^-$ generate groups of order $120$ with scalar subgroup $\{\pm1\}$, quotient of order $60$ with element orders $\{1,2,3,5\}$, consistent with $A_5$. (An older entry \texttt{Weil\_Z5\_*} in \texttt{checks.json} with ``order\_mod\_scalars'' $81$, $67$ is superseded by \texttt{weil5.json} and not used.) \emph{Reading (speculative):} the binary icosahedral factor in the cusp sector is formally reminiscent of ``lift invariant $-1$'' in the multi-$k$ language. This is a formal contact only: nothing here is a monodromy representation of a Hurwitz family, it does \textbf{not} meet the locked U1-E2 success criterion, it does not force Necessity, and it imports nothing into the residual core. Untested suggestion: whether the degree-$12$ braid-orbit permutation representation factors through a congruence quotient sharing constituents with $W_5^\pm$.}

%---------------------------------------------------------------------
\subsection{K3}\label{sec:rev8K3impl}
\begin{enumerate}[label=(\alph*),leftmargin=2em]
\item \emph{What is justified ($N_2$ only).} The Bridge-C lattice at $s=7$ enters K3 geometry through its even sublattice $N_2$: the $140$-label representation is $\rho_{N_2}\cong\rho_T^*$ with $T$ the transcendental lattice of an $N_2$-polarised K3 surface (Theorem~\ref{thm:weil8}, Corollary~\ref{cor:T8}; Category~A), such surfaces exist and have no $(-2)$-curves and no elliptic fibrations, and $\gamma^2$ lifts to an automorphism of entropy $2\log(6+\sqrt{35})$, with $\mathrm{Aut}(X)\cong D_\infty$ for very general $X$ (Theorems~\ref{thm:K3exist8}, \ref{thm:aut8}; Category~A (cited)). All K3 statements of this paper are for $N_2$ only, not for $N_1=(\Z^2,2A)$ (Remark~\ref{rem:open8}(e)).
\item \emph{Correction: $H^3(\mathrm{K3})=0$.} The cohomology of a K3 surface is concentrated in degrees $0,2,4$ (ranks $1,22,1$), so $H^3(\mathrm{K3};\Z)=0$. The corpus phrase ``resonant 3-form $H_3$ \dots\ as lift of K3 cohomology'' therefore cannot refer to a class on the K3 surface itself; it could only live on a product such as $\mathrm{K3}\times(\text{3-manifold})$, which the note would have to specify. This corrects a Category~B note, not a Category~A lock.
\item \emph{Correction: $2\cos(2\pi/539.9)$ is algebraic} (IC-7). As printed, $539.9=5399/10$ with $5399$ prime, so $\zeta=e^{2\pi i/539.9}=\e(10/5399)$ is a primitive $5399$-th root of unity and $2\cos(2\pi/539.9)$ is an algebraic integer of degree $\varphi(5399)/2=2699$. This contradicts the remark ``$\xi=2\cos(2\pi/539.9)$ is not algebraic as printed'' in \path{POLYNOMIAL_ROOTS_TARGETS_2026-08-17.md}. If $539.9$ is a rounded measured value, $\zeta$ is not exactly defined. Either way it is unrelated to K3 cohomology.
\item \emph{Not supported.} Any mock-theta, K3 elliptic genus or Mathieu-moonshine link to $C$ or $g_s$ (Category~B, not asserted: the Mathieu mock form has a polar term and grows exponentially radially, whereas $C\sim I/t$, and Theorem~\ref{thm:E} excludes scalar mock laws of the listed kinds); any $\Delta=\eta^{24}$/Yau--Zaslow reading of the Bernoulli numerators $691,3617,43867$ in the $S_e$ coefficients, which enter generically through Euler--Maclaurin (Category~A for the occurrence and mechanism, Category~B as a pattern). The Mukai-lattice embedding of a signature-$(1,1)$ lattice is true but has no specific content for Bridge~C.
\end{enumerate}
\catA{(a) as labelled in \S\ref{sec:rev8k3}. (b) standard topology. (c) exact: \texttt{checks.json}, keys \texttt{5399\_isprime}, \texttt{2cos(2pi*10/5399) degree} (PARI); $\varphi(5399)=5398$.}
